\documentclass[10pt,a4paper]{article}

\usepackage[T1]{fontenc}
\usepackage[utf8]{inputenc}
\usepackage[portuguese]{babel}
\usepackage[margin=2.1cm,top=2.3cm,bottom=2.3cm]{geometry}
\usepackage{newtxtext,newtxmath}
\usepackage{inconsolata}
\usepackage{xcolor}
\usepackage{fancyvrb}
\usepackage{tcolorbox}
\usepackage{titlesec}
\usepackage{enumitem}
\usepackage{fancyhdr}
\usepackage[colorlinks=true,linkcolor=azul,urlcolor=azul]{hyperref}

\definecolor{azul}{HTML}{2A4B8D}
\definecolor{ferrugem}{HTML}{A63A26}
\definecolor{fundo}{HTML}{F4F4F0}
\definecolor{borda}{HTML}{D8D8D1}
\definecolor{cinza}{HTML}{5A5F66}
\definecolor{verde}{HTML}{2E7D4F}

\tcbuselibrary{listings,skins,breakable}

% ---- caixa de comando ------------------------------------------------
% tcblisting, nao tcolorbox+ttfamily: sem verbatim de verdade, '_' vira
% subscrito, '#' some e '$PATH' entra em modo matematico.
\newtcblisting{comando}{
  enhanced, listing only, breakable, sharp corners,
  colback=fundo, colframe=borda, boxrule=0.4pt,
  borderline west={1.6pt}{0pt}{azul},
  left=6pt, right=6pt, top=4pt, bottom=4pt,
  before skip=5pt, after skip=2pt,
  listing options={
    basicstyle=\ttfamily\footnotesize,
    breaklines=true, breakindent=14pt,
    columns=fullflexible, keepspaces=true, showstringspaces=false,
  },
}

% ---- linha "o que isto testa" ---------------------------------------
\newcommand{\teste}[1]{%
  \par\vspace{1pt}\noindent
  {\color{verde}\footnotesize$\vartriangleright$}\hspace{3pt}%
  {\color{cinza}\footnotesize #1}\par\vspace{7pt}}

\newcommand{\aviso}[1]{%
  \par\vspace{1pt}\noindent
  {\color{ferrugem}\footnotesize$\blacktriangleright$}\hspace{3pt}%
  {\color{cinza}\footnotesize #1}\par\vspace{7pt}}

\titleformat{\section}{\normalfont\large\bfseries\color{azul}}{\thesection.}{0.6em}{}
\titleformat{\subsection}{\normalfont\normalsize\bfseries}{\thesubsection.}{0.6em}{}
\titlespacing{\section}{0pt}{16pt}{6pt}
\titlespacing{\subsection}{0pt}{11pt}{4pt}

\pagestyle{fancy}\fancyhf{}
\renewcommand{\headrulewidth}{0.4pt}
\fancyhead[L]{\footnotesize\color{cinza}PHASIS --- lista de comandos}
\fancyhead[R]{\footnotesize\color{cinza}\thepage}

\setlength{\parindent}{0pt}
\setlength{\parskip}{4pt}

\begin{document}

\begin{center}
{\LARGE\bfseries PHASIS --- lista de comandos}\\[3pt]
{\color{cinza}Como rodar o programa, e que teste cada comando permite fazer}\\[2pt]
{\footnotesize\color{cinza}commit \texttt{982f376} \quad$\cdot$\quad 22 de agosto de 2026}
\end{center}

\vspace{4pt}

{\color{cinza}\small
Legenda: {\color{verde}$\vartriangleright$} o que o comando permite testar
\quad $\cdot$ \quad {\color{ferrugem}$\blacktriangleright$} armadilha.
Todos os caminhos são relativos à raiz do repositório, exceto onde o
comando começa com \texttt{cd dipole}.}

% =====================================================================
\section{Ambiente}

Só o \texttt{dipole} com \texttt{-{}-use-F3 1} e os geradores em Python
precisam de biblioteca externa. O PHASIS compila e roda com \texttt{g++}
e a STL, sem mais nada.

\begin{comando}
export PATH=/home/rafael/micromamba/envs/dis/bin:$PATH
export LD_LIBRARY_PATH=/home/rafael/micromamba/envs/dis/lib:$LD_LIBRARY_PATH
\end{comando}
\teste{LHAPDF 6.5.6, \texttt{yadism} 0.13.11 e os conjuntos
\texttt{CT10nlo} e \texttt{NNPDF31\_nlo\_as\_0118} vivem nesse env
micromamba, não no Python do sistema.}

\begin{comando}
cd dipole && make env
\end{comando}
\teste{Diz se o LHAPDF foi detectado. Sem as duas linhas acima ele
reporta \texttt{ausente} e um binário já compilado falha com
\texttt{libLHAPDF.so: cannot open shared object file}.}

% =====================================================================
\section{Compilar}

\begin{comando}
make -j8
\end{comando}
\teste{Constrói \texttt{trace}, \texttt{alpha\_scan},
\texttt{redshift\_scan} e as oito fases de teste. Também há
\texttt{CMakeLists.txt}, mas o Makefile não exige \texttt{cmake} no
PATH.}

\begin{comando}
cd dipole && make -j8 && make build/sigma_nuN -j8
\end{comando}
\teste{O \texttt{dipole} tem build próprio. O acoplamento com o PHASIS é
\emph{por arquivo} --- as tabelas de seção de choque --- e não por link,
então os dois nunca precisam compilar juntos.}

% =====================================================================
\section{Suíte de aceitação}

\begin{comando}
make test
\end{comando}
\teste{247 verificações em 8 fases, $\sim$75~s no total. Deve terminar
com \texttt{0 falhas} em todas. É o primeiro comando a rodar depois de
qualquer mudança.}

\begin{comando}
./build/test_phase1        # 0,02 s -- quadratura, unidades, casos analíticos
./build/test_phase2        # 0,03 s -- relatividade geral, ponto de retorno, deflexão
./build/test_phase3        # 10 s   -- perfil com theta, rota de EDO acoplada
./build/test_phase4        #  6 s   -- cascata NC analítica, kernel, round-trip
./build/test_phase5        # 0,04 s -- casca tau = 1, protocolo de raízes
./build/test_phase6        # 0,35 s -- topologia do raio emitido, cone de escape
./build/test_phase7        # 0,06 s -- acoplamento GR x saturação, alpha_eff
./build/test_phase8        # 58 s   -- corrente neutra de ponta a ponta
\end{comando}
\teste{Rodar isolado quando se quer o detalhe. Cada um imprime linha a
linha o valor obtido, a referência e o desvio relativo.}
\aviso{As fases 7 e 8 leem as tabelas em \texttt{dipole/data/}. Se elas
não existirem, o teste \emph{falha} com a mensagem de como gerá-las ---
não pula em silêncio.}

\begin{comando}
cd dipole && make check
\end{comando}
\teste{Roda \texttt{test-quadrature}, \texttt{test-table} e
\texttt{validate-hera}. O primeiro compara $|\psi_T|^2$ e $|\psi_L|^2$
contra Kutak--Kwieciński eqs.~(12) e (13); deve dar desvio
$0{,}00\%$.}

% =====================================================================
\section{Gerar as tabelas do dipolo}

\subsection{Corrente carregada}

\begin{comando}
cd dipole
./build/sigma_nuN --model GBW --current CC --NE 300 --logEmin 3 --logEmax 14 \
    --nodesQ 16 --nodesX 16 --Ntabx 121 --NtabQ 121 \
    --use-F3 1 --pdf-set NNPDF31_nlo_as_0118
\end{comando}
\teste{$\sim$50~s. Produz \texttt{data/sigma\_nuN\_CC\_GBW.dat} com as
nove chaves de metadados e o hash do commit. Trocar \texttt{-{}-model
GBW} por \texttt{IIM} dá o bCGC.}

\subsection{Corrente neutra, total e diferencial}

\begin{comando}
./build/sigma_nuN --model GBW --current NC --NE 300 --logEmin 3 --logEmax 14 \
    --nodesQ 16 --nodesX 16 --nY 120 \
    --use-F3 1 --pdf-set NNPDF31_nlo_as_0118
\end{comando}
\teste{$\sim$1,5~min (o NC tem 4 canais de dipolo contra 2 do CC).
Produz \texttt{sigma\_nuN\_NC\_GBW.dat} \emph{e}
\texttt{dsigma\_dy\_NC\_GBW.dat}. O cabeçalho do segundo reporta, por
energia, que fração de $\sigma$ a integral em $y$ recupera --- deve
ficar entre 0,999 e 1,011.}

\aviso{O teste de física que vale fazer: $\sigma_{NC}/\sigma_{CC}$ tem de
convergir para $\approx 0{,}42$. Foi ele, e não nada interno ao código,
que pegou uma tabela NC com os acoplamentos errados que dava 2,49.}

\begin{comando}
python3 -c "
import numpy as np
cc=np.loadtxt('data/sigma_nuN_CC_GBW.dat'); nc=np.loadtxt('data/sigma_nuN_NC_GBW.dat')
E,r=cc[:,0],nc[:,2]/cc[:,2]
for v in (1e5,1e9,1e13):
    k=np.argmin(abs(E-v)); print('E=%9.2g  NC/CC=%.4f'%(E[k],r[k]))
print('previsto: (1.3127/4)*(M_Z/M_W)^2 = 0.4224')"
\end{comando}
\teste{Deve dar 0,358 / 0,405 / 0,412. A subida com a energia é o $xF_3$
do CC saindo de cena --- ele é de valência.}

\subsection{Convergência da quadratura}

\begin{comando}
for n in 16 48 128; do
  ./build/sigma_nuN --model GBW --current CC --NE 80 --logEmin 7 --logEmax 14 \
      --use-F3 0 --nodesQ $n --nodesX $n --quiet
  cp data/sigma_nuN_CC_GBW.dat /tmp/nos_$n.dat
done
\end{comando}
\teste{As três saídas devem concordar em $\sim 5\times10^{-7}$. Se
divergirem em $10^{-3}$ e a diferença cair como $N^{-1}$ em vez de
$N^{-4}$, há descontinuidade num extremo de quadratura --- foi
exatamente esse o sintoma do bug do extremo $y=1$, que custava um fator
670 em ruído.}

\begin{comando}
./build/sigma_nuN --model GBW --current CC --NE 80 --logEmin 7 --logEmax 14 \
    --use-F3 0 --Ntabx 241 --NtabQ 241 --quiet
\end{comando}
\teste{Dobrar a tabela de $F$ não deve mudar \emph{nada}. Se mudar, o
erro está na interpolação; se não mudar mas o resultado variar com
\texttt{-{}-nodesQ}, está na quadratura de $\sigma$. Separar as duas
causas é o que o par de comandos permite.}

% =====================================================================
\section{Validar o dipolo contra dados}

\begin{comando}
cd dipole
./build/validate_hera --model GBW --Nr 400 --Nz 200 \
    --Q2min 0.25 --Q2max 50 --xmax 0.01 --flavors udsc --mq 0.14
\end{comando}
\teste{Limite eletromagnético contra 239 pontos de H1+ZEUS combinados,
EPJ~C75 (2015) 580. Deve dar $\langle$modelo/dado$\rangle = 1{,}046$.
Escreve \texttt{data/hera\_validation.dat} com pull por ponto.}

\begin{comando}
./build/validate_hera --model IIM --iim-x0 0.69e-3 --Nr 400 --Nz 200
\end{comando}
\teste{O bCGC. Vale variar \texttt{-{}-iim-x0}: com
$x_0 = 0{,}00069\times 10^{-6}$ --- o valor que estava no código, com o
expoente da \emph{incerteza} lido como multiplicador --- a razão cai para
$0{,}10$ contra o próprio dado a que o modelo foi ajustado.}

\begin{comando}
./build/dipole --mode wavefunctions --current NC --flavor u --Q2 10
./build/dipole --mode integrand --current CC --flavor u --x 1e-5 --Q2 10
./build/dipole --mode sigma --current CC --flavor u --x 1e-5 --Q2 10
\end{comando}
\teste{Diagnóstico ponto a ponto. O modo \texttt{integrand} é o que
mostra o pico em $r \sim 1/\varepsilon$ --- 0,025~GeV$^{-1}$ em
$Q^2 = M_W^2$ --- que uma grade uniforme em $r$ perde inteiro.}

% =====================================================================
\section{A tabela colinear}

Serve de referência: as mesmas convenções e a mesma fórmula mestra do
dipolo, com funções de estrutura colineares em vez do modelo de dipolo.

\begin{comando}
python3 tools/build_collinear_sf.py
\end{comando}
\teste{$\sim$13~min. Tabula $F_2$, $F_L$ e $xF_3$ com yadism NLO sobre
NNPDF3.1 numa grade $121\times71$ em $(x,Q^2)$.}
\aviso{Fatia por $Q^2$ e grava cada fatia, então o RSS fica constante em
$\sim$800~MB e uma interrupção custa uma fatia. Pedir os 8591 pontos de
uma vez custaria 8,4~GB e travaria a máquina --- já aconteceu. Rodar de
novo retoma de onde parou.}

\begin{comando}
python3 tools/make_collinear_table.py --NE 300
\end{comando}
\teste{$\sim$5~s. Integra com a \emph{mesma} fórmula mestra do C++. Uma
quarta coluna diz que fração de $\sigma$ vem de $x$ abaixo do grid do
PDF: 0\% até $2\times10^{10}$~GeV, 21\% em $1{,}3\times10^{13}$. Acima
disso $\alpha_{\rm eff}$ colinear é propriedade da extrapolação, não
medida.}
\aviso{Confere com a literatura: $\sigma_{CC}(10^6~{\rm GeV}) \approx
6\times10^{-34}$~cm$^2$. Se sair fora disso por mais de 20\%, o cartão de
teoria ou o conjunto de PDF mudou.}

% =====================================================================
\section{Medir o acoplamento GR $\times$ saturação}

\begin{comando}
./build/alpha_scan --dipole dipole/data/sigma_nuN_CC_GBW.dat \
                   --colinear data/sigma_nuN_CC_collinear.dat \
                   --out data/alpha_scan_GBW.csv --nE 240
\end{comando}
\teste{Tabula $\alpha_{\rm eff} = d\ln\sigma/d\ln E$ e o teto
$3^{\alpha/2}$ para as duas tabelas. O teto tem de ser
\emph{não-crescente} em $E \geq 10^7$~GeV --- é a assinatura da
saturação, e é ela que distingue o efeito de um mero reajuste de PDF.}
\aviso{Chama \texttt{assert\_comparable} antes de qualquer conta. Passar
uma tabela de $\nu$ e uma de $\bar\nu$, ou uma NC e uma CC, faz lançar.
Vale tentar de propósito, para ver o guarda funcionar:
\texttt{-{}-colinear dipole/data/sigma\_nuN\_NC\_GBW.dat}.}

\begin{comando}
./build/redshift_scan --dipole dipole/data/sigma_nuN_CC_GBW.dat \
                      --colinear data/sigma_nuN_CC_collinear.dat \
                      --out data/redshift_scan_GBW.csv --nE 9
\end{comando}
\teste{Mede $R(b,E) = (\tau_{\rm full} - \tau_{\rm frozen})/\tau_{\rm
frozen}$ com a \emph{mesma} geometria; só o argumento de $\sigma$ muda.
$R$ é invariante sob $\sigma \to k\sigma$, então isola a forma e não a
normalização.}
\aviso{Geometria configurável: \texttt{-{}-rs -{}-rho0 -{}-r0 -{}-p
-{}-rin -{}-rout}. Vale testar \texttt{-{}-rin 2.0}: com $r_{\rm in}$
acima da esfera de fótons ($1{,}5\,r_s$), os raios quase-críticos
enrolam no vazio e o teto $3^{\alpha/2}$ deixa de ser atingível. O
default 1,2 é deliberado.}

% =====================================================================
\section{Traçar raios}

\begin{comando}
./build/trace
\end{comando}
\teste{Sem argumento, imprime todas as chaves de configuração
reconhecidas. É a documentação canônica do formato.}

\begin{comando}
./build/trace data/exemplo_terra.cfg saida_terra.csv
./build/trace data/exemplo_disco.cfg saida_disco.csv
./build/trace data/exemplo_buraco_negro.cfg saida_bn.csv
\end{comando}
\teste{Três casos prontos. O CSV traz 23 colunas:
\texttt{E\_inf\_GeV, b\_cm, i\_rad, psi\_t\_rad, r\_min\_cm, tau,
P\_surv, tau\_inbound, tau\_outbound, column\_g\_cm2, path\_cm,
theta\_min, theta\_max, n\_disk\_crossings, crosses\_matter, captured,
near\_critical, deflection\_rad, E\_loc\_max\_GeV, winding\_turns,
tolerance\_met, error, error\_msg}. Os dois ramos aparecem
separados --- \texttt{tau\_inbound} e \texttt{tau\_outbound} --- e só
coincidem em perfil esférico.}
\aviso{\texttt{exemplo\_buraco\_negro.cfg} é caso de \emph{estresse
numérico}, não cenário físico: põe matéria na esfera de fótons, o que
não tem órbita estável. Está no cabeçalho do próprio arquivo.}

\begin{comando}
# raio quase-critico: b logo acima de b_crit = 2.5980762 r_s
cat > /tmp/critico.cfg <<'EOF'
metric  = schwarzschild
r_s_cm  = 1.0e6
density    = powerlaw_halo
rho0_g_cm3 = 3.0
r0_cm      = 1.0e8
p          = 2.0
r_in_cm    = 1.2e6
r_out_cm   = 1.0e10
xsec       = powerlaw
sigma0_cm2 = 1e-33
E0_GeV     = 1e6
alpha      = 0.36
E_inf_GeV  = 1e8
b_min_cm   = 2.5980763e6
b_max_cm   = 3.0e6
n_b        = 20
deflection = 1
EOF
./build/trace /tmp/critico.cfg /tmp/critico.csv
\end{comando}
\teste{Varre o regime quase-crítico. $\tau$ vai de 199 a 1384 --- diverge
\emph{logaritmicamente} quando $b \to b_{\rm crit}^+$, e isso é físico:
o raio enrola na esfera de fótons e o comprimento próprio cresce mesmo.
A coluna \texttt{r\_min\_cm} tende a $1{,}5\,r_s = 1{,}5\times 10^6$.}
\aviso{O resumo reporta \texttt{degradados: 7} de 20 --- sete raios não
atingiram a tolerância. Não é falha silenciosa: a coluna
\texttt{tolerance\_met} marca quais. Apertar com
\texttt{rel\_tol = 1e-10} ou aumentar \texttt{max\_evals} resolve, ao
custo de tempo. Um resultado que não atingiu a tolerância pedida é
sempre sinalizado, nunca entregue como se tivesse atingido.}

\begin{comando}
# no mesmo arquivo, trocar as quatro linhas de xsec por:
#   xsec       = table
#   table_path = dipole/data/sigma_nuN_CC_GBW.dat
./build/trace /tmp/critico.cfg /tmp/critico_tabela.csv
\end{comando}
\teste{Troca a lei de potência pela tabela real. Com
\texttt{xsec\_current = total} é \emph{obrigatório} dar
\texttt{table\_path\_nc} ou \texttt{nc\_to\_cc\_ratio}: não há razão
NC/CC assumida por default, porque as tabelas do dipole são CC puras.}

% =====================================================================
\section{Plots}

\begin{comando}
python3 plots/literature.py
\end{comando}
\teste{Quatro figuras em \texttt{plots/}, PDF vetorial e PNG:
$\sigma(E)$ contra Gandhi et al.\ (via Formaggio \& Zeller, RMP~84
(2012) 1307, eqs.~90--91), $\alpha_{\rm eff}$ e o teto,
$\sigma_{NC}/\sigma_{CC}$, e o HERA.}
\aviso{As curvas da referência são desenhadas só na faixa em que ela se
declara válida ($10^{16}$ a $10^{21}$~eV). Estendê-las seria atribuir-lhe
uma afirmação que não faz, e justo onde a saturação muda a resposta.}

\begin{comando}
cd dipole && make plot-sigma-nuN plot-scan-x plot-hera-validation
\end{comando}
\teste{Alvos de plot do dipole. \texttt{make} sozinho lista os
disponíveis no início do arquivo.}

% =====================================================================
\section{O que ainda não funciona}

\begin{comando}
# NAO faca isto esperando resultado correto:
#   raio EMITIDO a distancia finita + perfil nao esferico
#   (density = flared_thin_disk, ou opts.force_ode)
\end{comando}
\aviso{A rota de EDO nunca recebeu a topologia de emissão. Para raio
emitido ela integra os dois ramos como se viesse do infinito: erra por
fator 7 num caso, lança exceção noutro, e devolve $\tau = 0$ em silêncio
no caso radial. Perfil esférico com raio vindo do infinito está
validado; o resto não. Idem para \texttt{transport\_cascade}.}

\begin{comando}
python3 -c "
import sys; sys.path.insert(0,'dipole/tests/oracle')
import dipole_oracle as o
print('oraculo  %.4e cm2' % o.sigma_cm2(1e9, NlogQ=8, Nlogx=8, Nr=100, Nz=100))"

cd dipole && ./build/sigma_nuN --model GBW --current CC \
    --NE 1 --logEmin 9 --logEmax 9 --use-F3 0 --quiet
\end{comando}
\aviso{Dá $7{,}628\times10^{-33}$ contra $7{,}352\times10^{-33}$ --- 3,8\%
de diferença. Mas a comparação hoje não isola nada, porque o oráculo
está atrás de \emph{duas} mudanças estruturais: a regra de densidade de
nós por década (a quadratura dele ainda usa número fixo) e o conserto do
extremo $y=1$. E não conhece NC. Ele existe para provar que uma mudança
no C++ fez o que se pretendia; enquanto estiver divergente, não serve
para isso. É o problema \textbf{P4} do relatório de estado.}

\vfill
{\footnotesize\color{cinza}
Gerado a partir de \texttt{docs/comandos.tex}. Os tempos são medidos
nesta máquina, com \texttt{OMP\_NUM\_THREADS} entre 4 e 6.
A máquina tem 15~GB de RAM e 2~GB de swap, sem \texttt{systemd-oomd}:
estourar a memória a trava em vez de matar o processo.}

\end{document}
